% =====================================================================
\section{Conclusion}\label{sec:conclusion}
% =====================================================================

\placeholder{Summary of the diode I--V results: $V_{th}$, $V_a$, $R_s$, $n$, leakage, effect of
light and of the junction shape (3--4 lines).}

\textbf{MIS capacitor.} The C--V curves at \SI{1}{MHz} show clearly the accumulation, depletion and
inversion regimes. The oxide thickness ($\approx\SI{45}{nm}$) is in agreement with the target of the
oxidation step, and the three methods give the same substrate doping, $N_A\approx\SI{1.3e16}{cm^{-3}}$.
However, the flat-band voltage ($\approx\SI{-4.2}{V}$) is far from the expected value
($\Phi_{MS}\approx\SI{-0.9}{V}$), which shows a large positive charge in the oxide
($N_s\approx\SI{1.5e12}{cm^{-2}}$). The hysteresis confirms that the oxide traps charges. The quality of
the oxide could probably be improved with a better cleaning before oxidation and an annealing step
after it. The measurements at low frequency and with light showed the role of the electrons
(minority carriers) in the inversion regime. The C--V curve of the diode gives a doping of the same
order of magnitude as the MIS capacitor.

\textbf{What this lab brought us.} \tocheck{to adapt} It was the first time we used a probe station and
an LCR meter on real devices. We saw how the measurement conditions (frequency, light, sweep direction)
change the curves, and that the values given directly by the software must be checked. Using several
methods to extract the same parameter also helped us to judge how reliable our results are.
